% ============================================================================
%  记号系统：preamble/math-notation.tex
%  ★ 全稿所有数学符号的唯一定义处。新增符号请写在这里并在附录符号表同步。
%  ★ 采用「\x 系列」命名，避免 \a \b \c 之类与已有宏冲突。
% ============================================================================

% --- 常用集合与空间 ---
\newcommand{\Rset}{\mathbb{R}}              % 实数集
\newcommand{\Nset}{\mathbb{N}}              % 自然数集
\newcommand{\Zset}{\mathbb{Z}}              % 整数集

% --- 运算符 ---
\DeclareMathOperator*{\argmin}{arg\,min}    % 带下划线的 argmin
\DeclareMathOperator*{\argmax}{arg\,max}
\DeclareMathOperator{\sign}{sgn}
\DeclareMathOperator{\tr}{tr}               % 矩阵迹

% --- 范数与括号 ---
\DeclarePairedDelimiter{\abs}{\lvert}{\rvert}
\DeclarePairedDelimiter{\norm}{\lVert}{\rVert}
\DeclarePairedDelimiter{\inner}{\langle}{\rangle}
\DeclarePairedDelimiterX{\set}[1]{\{}{\}}{#1}

% --- 偏导与微分 ---
\newcommand{\pd}[2]{\frac{\partial #1}{\partial #2}}
\newcommand{\dd}[1]{\mathrm{d}#1}

% --- 论文专属符号（示例，按实际模型增删）---
% 风险类
\newcommand{\pcrash}[0]{P_{\mathrm{crash}}}              % 坠机概率
\newcommand{\pcrashT}[2]{\pcrash\!\left(#1 \mid #2\right)} % 时变坠机概率 P(t | x, \theta)
\newcommand{\hazard}[0]{\lambda}                          % 危险率
\newcommand{\survival}[0]{S}                              % 生存函数
\newcommand{\exposure}[0]{\rho}                           % 暴露密度
\newcommand{\casualty}[0]{C_{\mathrm{fat}}}               % 致死代价
\newcommand{\propertyloss}[0]{C_{\mathrm{prop}}}          % 财产代价
\newcommand{\noisecost}[0]{C_{\mathrm{noise}}}            % 噪声代价
\newcommand{\riskcost}[0]{\mathcal{R}}                    % 综合风险代价
\newcommand{\riskfield}[0]{\mathcal{F}}                   % 风险张量场

% 规划类
\newcommand{\statespace}[0]{\mathcal{S}}    % 状态空间
\newcommand{\action}[0]{\mathcal{A}}        % 动作集
\newcommand{\nodeset}[0]{\mathcal{V}}       % 节点集
\newcommand{\edgeset}[0]{\mathcal{E}}       % 边集
\newcommand{\heur}[0]{h}                    % 启发函数
\newcommand{\gcost}[0]{g}                   % 累积代价
\newcommand{\tlabel}[0]{\tau}               % 时间标签

% --- 矩阵转置与向量 ---
\newcommand{\transpose}[0]{^{\mathsf{T}}}
\newcommand{\superscript}[1]{^{\mathrm{#1}}}
\newcommand{\subscript}[1]{_{\mathrm{#1}}}
